\documentclass{amsart}
% For dealing with references we use the comment environment
\usepackage{verbatim}
\newenvironment{reference}{\comment}{\endcomment}
%\newenvironment{reference}{}{}
\newenvironment{slogan}{\comment}{\endcomment}
\newenvironment{history}{\comment}{\endcomment}

% For commutative diagrams we use Xy-pic
\usepackage[all]{xy}

% We use 2cell for 2-commutative diagrams.
\xyoption{2cell}
\UseAllTwocells

% We use multicol for the list of chapters between chapters
\usepackage{multicol}

% This is generally recommended for better output
\usepackage{lmodern}
\usepackage[T1]{fontenc}

% For cross-file-references

% Package for hypertext links:
\usepackage{hyperref}

% For any local file, say "hello.tex" you want to link to please

% Theorem environments.
%
\theoremstyle{plain}
\newtheorem{theorem}[subsection]{Theorem}
\newtheorem{proposition}[subsection]{Proposition}
\newtheorem{lemma}[subsection]{Lemma}

\theoremstyle{definition}
\newtheorem{definition}[subsection]{Definition}
\newtheorem{example}[subsection]{Example}
\newtheorem{exercise}[subsection]{Exercise}
\newtheorem{situation}[subsection]{Situation}

\theoremstyle{remark}
\newtheorem{remark}[subsection]{Remark}
\newtheorem{remarks}[subsection]{Remarks}

\numberwithin{equation}{subsection}

% Macros
%
\def\lim{\mathop{\mathrm{lim}}\nolimits}
\def\colim{\mathop{\mathrm{colim}}\nolimits}
\def\Spec{\mathop{\mathrm{Spec}}}
\def\Hom{\mathop{\mathrm{Hom}}\nolimits}
\def\Ext{\mathop{\mathrm{Ext}}\nolimits}
\def\SheafHom{\mathop{\mathcal{H}\!\mathit{om}}\nolimits}
\def\SheafExt{\mathop{\mathcal{E}\!\mathit{xt}}\nolimits}
\def\Sch{\mathit{Sch}}
\def\Mor{\mathop{\mathrm{Mor}}\nolimits}
\def\Ob{\mathop{\mathrm{Ob}}\nolimits}
\def\Sh{\mathop{\mathit{Sh}}\nolimits}
\def\NL{\mathop{N\!L}\nolimits}
\def\CH{\mathop{\mathrm{CH}}\nolimits}
\def\proetale{{pro\text{-}\acute{e}tale}}
\def\etale{{\acute{e}tale}}
\def\QCoh{\mathit{QCoh}}
\def\Ker{\mathop{\mathrm{Ker}}}
\def\Im{\mathop{\mathrm{Im}}}
\def\Coker{\mathop{\mathrm{Coker}}}
\def\Coim{\mathop{\mathrm{Coim}}}

% Boxtimes
%
\DeclareMathSymbol{\boxtimes}{\mathbin}{AMSa}{"02}

%
% Macros for moduli stacks/spaces
%
\def\QCohstack{\mathcal{QC}\!\mathit{oh}}
\def\Cohstack{\mathcal{C}\!\mathit{oh}}
\def\Spacesstack{\mathcal{S}\!\mathit{paces}}
\def\Quotfunctor{\mathrm{Quot}}
\def\Hilbfunctor{\mathrm{Hilb}}
\def\Curvesstack{\mathcal{C}\!\mathit{urves}}
\def\Polarizedstack{\mathcal{P}\!\mathit{olarized}}
\def\Complexesstack{\mathcal{C}\!\mathit{omplexes}}
% \Pic is the operator that assigns to X its picard group, usage \Pic(X)
% \Picardstack_{X/B} denotes the Picard stack of X over B
% \Picardfunctor_{X/B} denotes the Picard functor of X over B
\def\Pic{\mathop{\mathrm{Pic}}\nolimits}
\def\Picardstack{\mathcal{P}\!\mathit{ic}}
\def\Picardfunctor{\mathrm{Pic}}
\def\Deformationcategory{\mathcal{D}\!\mathit{ef}}

\setlength{\emergencystretch}{3em}
\begin{document}
\title[Stacks Project, Tag 0AN5]{725 --- Formal spectrum maps of weakly admissible rings are representable exactly when taut}
\maketitle
\noindent Copyright (C) 2005--2025 Johan de Jong. Stacks Project authors. Source: \href{https://stacks.math.columbia.edu/tag/0AN5}{Tag 0AN5}.

\noindent Editorial difficulty note (not author text). Difficulty H2: The proof identifies pullbacks to infinitesimal affine thickenings through completed tensor products and proves that closure ideals form the required topology. The reverse direction collapses the completed tensor-product limit using tautness, a substantive formal representability criterion..

\noindent Editorial provenance note (not author text). History: excerpt from revision \texttt{a0444\allowbreak 6e57e\allowbreak c1fbc\allowbreak 252a8\allowbreak 71afc\allowbreak ec775\allowbreak 2fb28\allowbreak 07b14}, formal-spaces.tex, lines 4021--4096. Reference presentation was adapted; original-author context and core support blocks were retained or added. The mathematical wording is unchanged. Licensed under GNU FDL 1.2 or later; no invariant sections or cover texts. See ../licenses/COPYING. Context blocks reproduce author definitions and supporting results; they are not additional counted problems.

\begin{definition}
\label{definition-taut}
Let $\varphi : A \to B$ be a continuous map of linearly topologized rings.
We say $\varphi$ is {\it taut}\footnote{This is nonstandard notation.
The definition generalizes to modules, by saying a linearly topologized
$A$-module $M$ is $A$-taut if for every open ideal $I \subset A$ the closure
of $IM$ in $M$ is open and these closures form a fundamental system of
neighbourhoods of $0$ in $M$.}
if for every open ideal $I \subset A$ the closure of the ideal $\varphi(I)B$
is open and these closures form a fundamental system of open ideals.
\end{definition}

\begin{definition}
\label{definition-affine-formal-algebraic-space}
Let $S$ be a scheme. We say a sheaf $X$ on $(\Sch/S)_{fppf}$ is an
{\it affine formal algebraic space} if there exist
\begin{enumerate}
\item a directed set $\Lambda$,
\item a system $(X_\lambda, f_{\lambda \mu})$ over $\Lambda$
in $(\Sch/S)_{fppf}$ where
\begin{enumerate}
\item each $X_\lambda$ is affine,
\item each $f_{\lambda \mu} : X_\lambda \to X_\mu$ is a thickening,
\end{enumerate}
\end{enumerate}
such that
$$
X \cong \colim_{\lambda \in \Lambda} X_\lambda
$$
as fppf sheaves and $X$ satisfies a set theoretic condition
(see Remark \href{https://stacks.math.columbia.edu/tag/0AIS}{remark set theoretic (Tag 0AIS)}). A
{\it morphism of affine formal algebraic spaces}
over $S$ is a map of sheaves.
\end{definition}

\begin{lemma}
\label{lemma-mcquillan-affine-formal-algebraic-space}
Let $S$ be a scheme. Let $X$ be an fppf sheaf on $(\Sch/S)_{fppf}$
which satisfies the set theoretic condition of
Remark \href{https://stacks.math.columbia.edu/tag/0AIS}{remark set theoretic (Tag 0AIS)}.
The following are equivalent:
\begin{enumerate}
\item there exists a weakly admissible topological ring $A$ over $S$
(see Remark \href{https://stacks.math.columbia.edu/tag/0AI3}{remark mcquillan (Tag 0AI3)}) such that
$X = \colim_{I \subset A\text{ weak ideal of definition}} \Spec(A/I)$,
\item $X$ is an affine formal algebraic space and
there exists an $S$-algebra $A$ and a map $X \to \Spec(A)$
such that for a closed immersion $T \to X$ with $T$ an affine scheme
the composition $T \to \Spec(A)$ is a closed immersion,
\item $X$ is an affine formal algebraic space and
there exists an $S$-algebra $A$ and a map $X \to \Spec(A)$
such that for a closed immersion $T \to X$ with $T$ a scheme
the composition $T \to \Spec(A)$ is a closed immersion,
\item $X$ is an affine formal algebraic space and
for some choice of $X = \colim X_\lambda$ as in
Definition \ref{definition-affine-formal-algebraic-space}
the projections $\lim \Gamma(X_\lambda, \mathcal{O}_{X_\lambda})
\to \Gamma(X_\lambda, \mathcal{O}_{X_\lambda})$ are surjective,
\item $X$ is an affine formal algebraic space and for any choice
of $X = \colim X_\lambda$ as in
Definition \ref{definition-affine-formal-algebraic-space}
the projections $\lim \Gamma(X_\lambda, \mathcal{O}_{X_\lambda})
\to \Gamma(X_\lambda, \mathcal{O}_{X_\lambda})$ are surjective.
\end{enumerate}
Moreover, the weakly admissible topological ring is
$A = \lim \Gamma(X_\lambda, \mathcal{O}_{X_\lambda})$
endowed with its limit topology and the weak ideals of definition
classify exactly the morphisms $T \to X$ which are representable
and thickenings.
\end{lemma}

\begin{proof}
It is clear that (5) implies (4).

\medskip\noindent
Assume (4) for $X = \colim X_\lambda$ as in
Definition \ref{definition-affine-formal-algebraic-space}.
Set $A = \lim \Gamma(X_\lambda, \mathcal{O}_{X_\lambda})$.
Let $T \to X$ be a closed immersion with $T$ a scheme
(note that $T \to X$ is representable by
Lemma \href{https://stacks.math.columbia.edu/tag/0AI8}{lemma diagonal affine formal algebraic space (Tag 0AI8)}).
Since $X_\lambda \to X$ is a thickening, so is
$X_\lambda \times_X T \to T$. On the other hand,
$X_\lambda \times_X T \to X_\lambda$ is a closed immersion,
hence $X_\lambda \times_X T$ is affine. Hence $T$ is affine
by Limits, Proposition \href{https://stacks.math.columbia.edu/tag/05YU}{proposition affine (Tag 05YU)}.
Then $T \to X$ factors through $X_\lambda$ for some $\lambda$
by Lemma \href{https://stacks.math.columbia.edu/tag/0AIA}{lemma factor through thickening (Tag 0AIA)}.
Thus $A \to \Gamma(X_\lambda, \mathcal{O}) \to \Gamma(T, \mathcal{O})$
is surjective. In this way we see that (3) holds.

\medskip\noindent
It is clear that (3) implies (2).

\medskip\noindent
Assume (2) for $A$ and $X \to \Spec(A)$. Write $X = \colim X_\lambda$
as in Definition \ref{definition-affine-formal-algebraic-space}.
Then $A_\lambda = \Gamma(X_\lambda, \mathcal{O})$ is a quotient
of $A$ by assumption (2). Hence $A^\wedge = \lim A_\lambda$
is a complete topological ring, see discussion in
More on Algebra, Section \href{https://stacks.math.columbia.edu/tag/07E7}{section topological ring (Tag 07E7)}.
The maps $A^\wedge \to A_\lambda$ are surjective as $A \to A_\lambda$ is.
We claim that for any $\lambda$ the kernel $I_\lambda \subset A^\wedge$ of
$A^\wedge \to A_\lambda$ is a weak ideal of definition.
Namely, it is open by definition of the limit topology.
If $f \in I_\lambda$, then for any $\mu \in \Lambda$
the image of $f$ in $A_\mu$ is zero in all the residue fields
of the points of $X_\mu$. Hence it is a nilpotent element
of $A_\mu$. Hence some power $f^n \in I_\mu$. Thus $f^n \to 0$
as $n \to 0$. Thus $A^\wedge$ is weakly admissible.
Finally, suppose that $I \subset A^\wedge$ is a weak ideal
of definition. Then $I \subset A^\wedge$ is open and hence there exists
some $\lambda$ such that $I \supset I_\lambda$. Thus we obtain a morphism
$\Spec(A^\wedge/I) \to \Spec(A_\lambda) \to X$.
Then it follows that $X = \colim \Spec(A^\wedge/I)$ where now
the colimit is over all weak ideals of definition.
Thus (1) holds.

\medskip\noindent
Assume (1). In this case it is clear that $X$ is an affine formal
algebraic space. Let $X = \colim X_\lambda$ be any presentation as in
Definition \ref{definition-affine-formal-algebraic-space}.
For each $\lambda$ we can find a weak ideal of definition
$I \subset A$ such that $X_\lambda \to X$ factors through
$\Spec(A/I) \to X$, see Lemma \href{https://stacks.math.columbia.edu/tag/0AIA}{lemma factor through thickening (Tag 0AIA)}.
Then $X_\lambda = \Spec(A/I_\lambda)$ with $I \subset I_\lambda$.
Conversely, for any weak ideal of definition $I \subset A$
the morphism $\Spec(A/I) \to X$ factors through $X_\lambda$
for some $\lambda$, i.e., $I_\lambda \subset I$.
It follows that each $I_\lambda$ is a weak ideal of definition
and that they form a cofinal subset of the set of weak ideals
of definition. Hence $A = \lim A/I = \lim A/I_\lambda$
and we see that (5) is true and moreover that
$A = \lim \Gamma(X_\lambda, \mathcal{O}_{X_\lambda})$.
\end{proof}


\begin{lemma}
\label{lemma-representable-affine}
Let $S$ be a scheme. Let $\varphi : A \to B$ be a continuous map of
weakly admissible topological rings over $S$. The following
are equivalent
\begin{enumerate}
\item $\text{Spf}(\varphi) : \text{Spf}(B) \to \text{Spf}(A)$
is representable by algebraic spaces,
\item $\text{Spf}(\varphi) : \text{Spf}(B) \to \text{Spf}(A)$
is representable (by schemes),
\item $\varphi$ is taut, see Definition \ref{definition-taut}.
\end{enumerate}
\end{lemma}

\begin{proof}
Parts (1) and (2) are equivalent by
Lemma \href{https://stacks.math.columbia.edu/tag/0AKN}{lemma affine representable by algebraic spaces (Tag 0AKN)}.

\medskip\noindent
Assume the equivalent conditions (1) and (2) hold.
If $I \subset A$ is a weak ideal of definition, then
$\Spec(A/I) \to \text{Spf}(A)$ is representable and a thickening
(this is clear from the construction of the formal spectrum
but it also follows from
Lemma \ref{lemma-mcquillan-affine-formal-algebraic-space}).
Then $\Spec(A/I) \times_{\text{Spf}(A)} \text{Spf}(B) \to \text{Spf}(B)$
is representable and a thickening as a base change.
Hence by
Lemma \ref{lemma-mcquillan-affine-formal-algebraic-space}
there is a weak ideal of definition $J(I) \subset B$ such that
$\Spec(A/I) \times_{\text{Spf}(A)} \text{Spf}(B) = \Spec(B/J(I))$
as subfunctors of $\text{Spf}(B)$. We obtain a cartesian diagram
$$
\xymatrix{
\Spec(B/J(I)) \ar[d] \ar[r] & \Spec(A/I) \ar[d] \\
\text{Spf}(B) \ar[r] & \text{Spf}(A)
}
$$
By Lemma \href{https://stacks.math.columbia.edu/tag/0AN3}{lemma fibre product affines over separated (Tag 0AN3)}
we see that $B/J(I) = B \widehat{\otimes}_A A/I$.
It follows that $J(I)$ is the closure of the ideal $\varphi(I)B$, see
Lemma \href{https://stacks.math.columbia.edu/tag/0AMZ}{lemma closure image ideal (Tag 0AMZ)}.
Since $\text{Spf}(A) = \colim \Spec(A/I)$ with $I$ as above,
we find that $\text{Spf}(B) = \colim \Spec(B/J(I))$.
Thus the ideals $J(I)$ form a fundamental system of weak
ideals of definition (see
Lemma \ref{lemma-mcquillan-affine-formal-algebraic-space}).
Hence (3) holds.

\medskip\noindent
Assume (3) holds. We are essentially just going to reverse the
arguments given in the previous paragraph.
Let $I \subset A$ be a weak ideal of definition.
By Lemma \href{https://stacks.math.columbia.edu/tag/0AN3}{lemma fibre product affines over separated (Tag 0AN3)}
we get a cartesian diagram
$$
\xymatrix{
\text{Spf}(B \widehat{\otimes}_A A/I) \ar[d] \ar[r] & \Spec(A/I) \ar[d] \\
\text{Spf}(B) \ar[r] & \text{Spf}(A)
}
$$
If $J(I)$ is the closure of $IB$, then $J(I)$ is open in $B$
by tautness of $\varphi$. Hence if $J$ is open in $B$ and $J \subset J(B)$,
then $B/J \otimes_A A/I = B/(IB + J) = B/J(I)$ because
$J(I) = \bigcap_{J \subset B\text{ open}} (IB + J)$ by Lemma \href{https://stacks.math.columbia.edu/tag/0AMS}{lemma closed (Tag 0AMS)}.
Hence the limit defining the completed tensor product collapses to give
$B \widehat{\otimes}_A A/I = B/J(I)$.
Thus $\text{Spf}(B \widehat{\otimes}_A A/I) = \Spec(B/J(I))$.
This proves that $\text{Spf}(B) \times_{\text{Spf}(A)} \Spec(A/I)$
is representable for every weak ideal of definition $I \subset A$.
Since every morphism $T \to \text{Spf}(A)$ with $T$ quasi-compact
factors through $\Spec(A/I)$ for some weak ideal of definition $I$
(Lemma \href{https://stacks.math.columbia.edu/tag/0AIA}{lemma factor through thickening (Tag 0AIA)})
we conclude that $\text{Spf}(\varphi)$ is representable, i.e.,
(2) holds. This finishes the proof.
\end{proof}
\end{document}
